\documentclass[a4paper,twoside]{article}


% --- RUIMTE EN LAYOUT ---
\usepackage{setspace}

\usepackage{geometry} %size regulating
\geometry{
 a4paper,
 left=25mm,
 right=25mm,
 top=20mm,
 bottom=20mm,
 }

\usepackage{parskip} % for better spacing between paragraphs
\usepackage{flushend} % Alleen nodig als je laatste pagina van artikel wilt vullen met tekst, zodat de kolommen even lang zijn. Anders kan je dit beter weglaten.

% --- WISKUNDE ENSYMBOLEN ---
\usepackage{amsmath, amssymb, amsfonts, amsthm} % For nicer math format
\usepackage{mathrsfs} % For curly arrr
\usepackage{bm} % For bold symbols eg \bm{\nabla}
\usepackage{esint} % For bolletje in integrals
\usepackage{dsfont}
\usepackage{siunitx} % Voeg deze expliciet toe voor die graden en eenheden!
\usepackage{braket} % For bra-ket notation
\usepackage{mathtools}
\usepackage{url}
\allowdisplaybreaks

% --- OVERIGE ---

\usepackage{graphicx} % Required for inserting images
\usepackage[hidelinks]{hyperref} % geen kaders om linken
\usepackage{subcaption}
%\usepackage[dutch]{babel}
\usepackage[T1]{fontenc}
\usepackage[utf8]{inputenc}
\usepackage{lmodern}
\usepackage{multicol}


% --- Tikz packages ---
\usepackage{tikz}
\usetikzlibrary{angles,quotes, babel} % for pic (angle labels)
\usetikzlibrary{arrows.meta,decorations.markings,calc, decorations.pathmorphing, decorations.pathreplacing, positioning, shapes, shadings, patterns}
\usepackage{xcolor}
\usepackage{tikz-3dplot}
\usepackage[siunitx]{circuitikz}

\usepackage{etoolbox} % ifthen
\usepackage[outline]{contour} % glow around text
\usepackage{xfp} % higher precision (16 digits?)
\contourlength{1.1pt}



% --- THEOREM STIJLEN MET EXTRA RUIMTE ---

\theoremstyle{definition}
\newtheorem{postulate}{Postulaat}
\newtheorem{definition}{Definitie}
\newtheorem{derivation}{Afleiding}
\newtheorem{proposition}{Propositie}
\newtheorem{example}{Voorbeeld}

\theoremstyle{plain}
\newtheorem{theorem}{Theorem}


% --- EXTRA COMMANDO'S ---

\newcommand{\dbar}{{\mkern+3mu\mathchar'26\mkern-12mu d}}

% --- Titel en Introductie --- 

\title{Vectoren en Functie Ruimten}
\author{T. Cools}
\date{2026-2027}
\raggedbottom

% Toelaten dat pagina-overgangen plaatsvinden tussen opeenvolgende (lege) titels
\makeatletter
\let\old@startsection\@startsection
\renewcommand{\@startsection}[6]{%
  \@nobreakfalse
  \old@startsection{#1}{#2}{#3}{#4}{#5}{#6}%
}
\makeatother

\begin{document}
\maketitle

\begin{abstract}
   Hier noteer ik de delen van de cursus die ik nog niet helemaal begrijp of moeilijk vind.
\end{abstract}

\section{ELEMENTARY ALGEBRAIC STRUCTURES}
\subsection{Sets, maps and relations}
\subsubsection{Sets}
\subsubsection{Maps}
\subsubsection{Set cardinality revisited}
\begin{itemize}
    \item cardinality of the continuum??
    \item Power sets of (in)finite sets \cite{enwiki:power_set}
\end{itemize}
\subsubsection{Relations}
\begin{itemize}
    \item Example 1.2 (oefening!)
    \item partial order relation: antisymmetric??
    \item Example 1.4 (oefening!)
\end{itemize}
\subsection{Groups, rings and fields}
\subsubsection{Groups}
\subsubsection{Homomorphisms, isomorphisms and automorphisms}
\subsubsection{Group actions}
\subsubsection{Kernels, normal subgroups and quotient groups}
\subsubsection{Rings and fields}
\subsection{Vector spaces}
\subsubsection{Definitions and examples}
\subsubsection{Vector space homomorphisms, endomorphisms and isomorphisms}
\subsubsection{Linear combination, span and completeness}
\subsubsection{Linear independence and basis}
\subsubsection{Index notation and Einstein summation convention}
\subsubsection{Combining and manipulating spaces and subspaces}
\subsection{Algebras}

\section{LINEAR MAPS AND MATRICES}
\subsection{Linear maps reviewed}
\subsection{Kernel, image and rank-nullity theorem}
\subsection{Matrices and determinants}
\subsubsection{Matrix representation of linear maps}
\subsubsection{Matrix properties and manipulations}
\subsection{Computation and complexity of matrix multiplication}
\subsubsection{Trace and determinant}
\subsubsection{Application: integration measures and Jacobians}
\subsubsection{Matrix inverse}
\subsection{General linear group and basis transforms}
\subsubsection{Matrix groups}
\subsubsection{Basis transforms}
\subsection{Functionals and dual spaces}
\subsubsection{Dual spaces}
\subsubsection{Basis transforms and the contragradient representation}
\subsubsection{Dual linear maps and the transpose}
\subsection{Systems of linear equations}
\subsubsection{Gaussian elimination and LU decomposition}

\section{LINEAR OPERATORS AND EIGENVALUES}
\subsection{Powers and polynomials}
\subsubsection{Powers and polynomials of operators}
\subsubsection{Projection operators}
\subsubsection{Annihilating polynomials}
\subsection{Eigenvalues, eigenvectors and diagonalisation}
\subsubsection{Eigenvalues, eigenvectors and eigenspaces}
\subsubsection{Characteristic polynomial and Cayley-Hamilton theorem}
\subsubsection{Diagonalisation and spectral decomposition}
\subsubsection{Diagonalisation of commuting operators}
\subsection{Generalised eigenspaces and Jordan canonical form}
\subsubsection{Invariant subspaces}
\subsubsection{Generalised eigenspaces}
\subsubsection{Jordan canonical form}
\subsubsection{Related eigenvalue problems}
\subsection{Functions of linear operators}
\subsubsection{Matrix exponential}
\subsection{Application: dynamical systems}
\subsubsection{Recurrence relations}
\subsubsection{Initial value problems}

\section{NORMS AND DISTANCES}
\subsection{Normed vector spaces}
\subsubsection{Hölder norms}
\subsubsection{Interlude: calculus in metric spaces}
\subsubsection{Convergence and continuity in normed vector spaces}
\subsubsection{Equivalence of norms}
\subsection{Banach spaces}
\subsection{Norms for linear maps}
\subsubsection{Continuity and subordinate norms}
\subsubsection{Matrix norms}
\subsubsection{Spectral radius and Gelfand formula}
\subsection{Applications}
\subsubsection{Sensitivity of linear systems}
\subsubsection{Functions of operators revisited}

\section{INNER PRODUCTS AND ORTHOGONALITY}
\subsection{Prologue: bilinear and sesquilinear forms}
\subsection{Inner product spaces}
\subsubsection{Cauchy-Schwarz inequality and its consequences}
\subsubsection{Examples of infinite-dimensional inner product spaces}
\subsubsection{Parallelogram law and polarisation identity}
\subsection{Orthogonality and unitarity}
\subsubsection{Orthogonality and orthonormality}
\subsubsection{Orthogonal complements and orthogonal projections}
\subsubsection{Orthonormal basis for Hilbert spaces}
\subsubsection{Gram-Schmidt orthonormalisation and the QR decomposition}
\subsection{Linear functionals and the duality of Hilbert space}
\subsubsection{Finite-dimensional case}
\subsubsection{Riesz representation theorem}
\subsection{Bounded linear maps in Hilbert spaces}
\subsubsection{Preliminaries}
\subsubsection{Adjoint of a linear map}
\subsubsection{Self-adjoint operators}
\subsubsection{Isometric and unitary maps}
\subsubsection{Normal operators}
\subsubsection{Hilbert-Schmidt inner product}
\subsection{Application: least squares solutions}

\section{MULTILINEAR ALGEBRA}
\subsection{Tensor product of vector spaces}
\subsubsection{Bilinear maps revisited}
\subsubsection{Formal definition}
\subsubsection{Tensor product of linear maps}
\subsubsection{Multilinear maps and general tensors}
\subsubsection{Dual spaces and linear maps revisited}
\subsubsection{Index notation and tensor operations}
\subsection{Tensors in inner product spaces}
\subsubsection{Inner product in a tensor space}
\subsubsection{Raising and lowering indices}
\subsubsection{Tensor products of infinite dimensional Hilbert spaces}
\subsection{Tensor algebra and exterior products}
\subsubsection{Tensor algebra}
\subsubsection{Symmetry under index permutation}
\subsubsection{Wedge products, exterior algebra and alternating forms}

\section{UNITARY SIMILARITY AND UNITARY EQUIVALENCE}
\subsection{Unitary and orthogonal groups}
\subsection{Elementary unitary transformations}
\subsubsection{Permutation matrices}
\subsubsection{Discrete Fourier transformation}
\subsection{Schur decomposition and power iteration}
\subsubsection{Schur decomposition}
\subsubsection{Normal matrices revisited}
\subsubsection{Power method and subspace iteration}
\subsection{Singular value and polar decomposition}
\subsubsection{Unitary equivalence}
\subsubsection{Singular value decomposition}
\subsubsection{Rank, norm, condition number}
\subsubsection{Least squares and pseudo-inverse}
\subsubsection{Low rank approximations}
\subsubsection{Polar decomposition}

\section{FUNCTION SPACES}
\subsection{Function spaces and norms}
\subsubsection{Introduction}


\bibliographystyle{plain}
\bibliography{/Users/tristancools/Documents/Overleaf-Latex-Typst/References/references}

\end{document}